\documentclass{article}
% ============================================================
% Bibliography and citations
% ============================================================

\usepackage{natbib}
\usepackage{csquotes}


% ============================================================
% Mathematics
% ============================================================

\usepackage{
    amsmath,
    amssymb,
    amsthm,
    amsfonts,
    mathrsfs,
    dsfont,
    bm,
    textgreek,
    bbm
}


% ============================================================
% Fonts
% ============================================================

\usepackage{newpxtext}
\usepackage{eulervm}


% ============================================================
% Page layout and document structure
% ============================================================

\usepackage[margin=1in, footskip=0.3in]{geometry}
\usepackage{setspace}
\usepackage[bf,small]{titlesec}
\usepackage{etoc}
\usepackage{chngcntr}
\usepackage{changepage}
\usepackage{afterpage}


% ============================================================
% Figures, captions, and graphics
% ============================================================

\usepackage{graphicx}
\usepackage[bf,skip=8pt]{caption}
\usepackage{subcaption}
\usepackage{tikz}

\graphicspath{{Figures/}}


% ============================================================
% Tables
% ============================================================

\usepackage{booktabs}
\usepackage{tabularx}
\usepackage{dcolumn}
\usepackage{makecell}
\usepackage{multirow}
\usepackage{longtable}
\usepackage[para]{threeparttable}
\usepackage{colortbl}
\usepackage{adjustbox}


% ============================================================
% Landscape and rotated content
% ============================================================

\usepackage{lscape}
\usepackage{pdflscape}
\usepackage{rotating}


% ============================================================
% Lists, comments, notes, and text utilities
% ============================================================

\usepackage{enumerate}
\usepackage{comment}
\usepackage{verbatim}
\usepackage{soul}
\usepackage[multiple,stable]{footmisc}
\usepackage[colorinlistoftodos]{todonotes}


% ============================================================
% Code listings
% ============================================================

\usepackage{listings}


% ============================================================
% Colors
% ============================================================

\usepackage{xcolor}

% General-purpose colors
\definecolor{blue}{rgb}{0,0.08,0.5}
\definecolor{red}{rgb}{0.6,0,0}
\definecolor{green}{rgb}{0,0.5,0.05}

% Code-listing colors
\definecolor{fgreen}{HTML}{698B27}
\definecolor{lavender}{HTML}{E6E6FA}
\definecolor{darkpink}{HTML}{C71585}
\definecolor{darkyellow}{HTML}{B8860B}


% ============================================================
% Hyperlinks and PDF metadata
% ============================================================

\usepackage[
    bookmarks,
    pdftitle={},
    pdfauthor={Banks},
    colorlinks=true,
    linkcolor=black,
    citecolor=black,
    urlcolor=black
]{hyperref}


% ============================================================
% Theorem environments
% ============================================================

\newtheorem{prediction}{Prediction}
\newtheorem{assumption}{Assumption}
\newtheorem{proposition}{Proposition}
\newtheorem{corollary}{Corollary}
\newtheorem{definition}{Definition}
\newtheorem{lemma}{Lemma}

% Proposition numbering resets by section
\newtheorem{prop}{Proposition}[section]


% ============================================================
% Mathematical commands
% ============================================================

% Formatting
\newcommand{\no}{\noindent}
\newcommand{\mb}{\mathbf}

% Expectation, probability, and covariance
\newcommand{\E}{\mathbb{E}}
\newcommand{\R}{\mathbb{R}}
\newcommand{\pr}{\mathrm{Pr}}
\newcommand{\cov}{\mathrm{Cov}}

% Support
\newcommand{\supp}[1]{\mathrm{Supp}\left(#1\right)}

% Partial derivatives
\newcommand{\pd}[2]{\frac{\partial #1}{\partial #2}}
\newcommand{\pdd}[2]{\frac{\partial^2 #1}{\partial #2^2}}

% Indicator function
\newcommand{\ind}[1]{\mathbbm{1}\{#1\}}

% Vertical spacing helper for large fractions
\newcommand{\vh}{
    \vphantom{
        \frac{
            \frac{\pd{a}{a}}{\pd{a}{a}}
        }{
            \frac{\pd{a}{a}}{\pd{a}{a}}
        }
    }
}


% ============================================================
% Custom environments
% ============================================================

% Wide equation extending 2 cm into each margin
\newenvironment{widerequation}
{
    \begin{adjustwidth}{-2cm}{-2cm}
    \begin{equation*}
}
{
    \end{equation*}
    \end{adjustwidth}
}

% Solution environment
\newenvironment{solution}
{
    \par\noindent\textbf{\textit{Solution.}}
}
{
    \hfill\rule{1.1ex}{1.1ex}
}


% ============================================================
% Code-listing style
% ============================================================

\lstset{
    basicstyle=\ttfamily\small\color{fgreen},
    keywordstyle=\color{darkpink},
    commentstyle=\color{gray},
    stringstyle=\color{darkyellow},
    showstringspaces=false,
    breaklines=true,
    frame=single,
    backgroundcolor=\color{lavender}
}


% ============================================================
% Spacing and section formatting
% ============================================================

% ECMA-style line spacing
\linespread{1.5}

\titlespacing*{\paragraph}
    {0pt}
    {.25\baselineskip}
    {1em}

\titlespacing*{\section}
    {0pt}
    {.6\baselineskip}
    {.3\baselineskip}

\titlespacing*{\subsection}
    {0pt}
    {.6\baselineskip}
    {.3\baselineskip}

\titlespacing*{\subsubsection}
    {0pt}
    {.6\baselineskip}
    {.3\baselineskip}

% \setlength{\parskip}{5pt}

\setlength{\footnotesep}{.2in}


% ============================================================
% Subsubsection formatting
% ============================================================

\titleformat{\subsubsection}[runin]
    {\normalfont\itshape\bfseries}
    {\thesubsubsection}
    {1em}
    {}
    [.]


% ============================================================
% Date and title-page utilities
% ============================================================

\newcommand{\monthyear}{%
    \ifcase\month\or
        January\or
        February\or
        March\or
        April\or
        May\or
        June\or
        July\or
        August\or
        September\or
        October\or
        November\or
        December
    \fi
    \ \number\year
}


\begin{document}

    \begin{center}
        \scalebox{1.3}{\textbf{Economics 880 -- PSet 3}}\\
        Aaron Banks \& Arthur Rodrigues  
    \end{center}

\section{Model Comparison}

\begin{table}[htbp]
    \centering
    \caption{Stationary Equilibrium: Model Comparison}
    \label{tab:equilibrium_comparison}
    \begin{tabular}{lccc}
        \hline\hline
        Variable & Standard & TV1 ($\alpha=1$) & TV1 ($\alpha=2$) \\
        \hline
        Price Level                  & 0.7390   & 0.6908   & 0.7195   \\
        Mass of Incumbents           & 8.3207   & 9.5489   & 8.3505   \\
        Mass of Entrants             & 2.6386   & 4.2232   & 3.5112   \\
        Mass of Exits                & 2.6386   & 4.2232   & 3.5112   \\
        \hline
        Aggregate Labor              & 179.8332 & 188.8909 & 182.6240 \\
        \quad Labor of Incumbents    & 142.6260 & 139.5075 & 136.6501 \\
        \quad Labor of Entrants      & 37.2072  & 49.3834  & 45.9739  \\
        Fraction of Labor in Entrants & 0.2069  & 0.2614   & 0.2517   \\
        \hline\hline
    \end{tabular}
\end{table}

The TV1 specifications generate lower equilibrium prices, greater firm turnover,
and higher aggregate employment relative to the standard model. These
differences can be understood through three mechanisms.

First, introducing action-specific shocks increases the ex-ante value of
operating a firm, since firms choose whether to exit or continue after
observing these shocks. This makes entry more attractive at a given output
price. Consequently, the equilibrium price must fall to restore the
free-entry condition. This effect is stronger when $\alpha$ is smaller,
corresponding to greater variance in the action-specific shocks.

Second, the TV1 shocks weaken the relationship between productivity and
exit. In the standard model, all firms in the lowest productivity state
exit, while firms in higher productivity states survive. Under TV1,
some low-productivity firms survive while some higher-productivity firms
exit. In our calibration, this increases firm turnover and reduces average
firm lifetimes. Since entry must equal exit in a stationary equilibrium,
the higher rate of exit is accompanied by greater entry. As a result,
the fraction of aggregate employment accounted for by entrants increases.

Finally, aggregate employment increases despite the lower equilibrium
price, which reduces labor demand at the individual-firm level. This
reflects adjustments along both the extensive margin of firm entry and
the household labor supply margin. In particular, aggregate operating
profits decline in the TV1 specifications. Since household labor supply
satisfies
\[
    N^s = \frac{1}{A} - \Pi,
\]
lower profit income induces households to supply more labor, supporting
higher equilibrium employment.

As $\alpha$ increases, the variance of the action-specific shocks
decreases, and the TV1 model approaches the deterministic specification.
Accordingly, the equilibrium price, aggregate employment, and firm
turnover generally move toward their standard-model values. However,
individual equilibrium moments need not change monotonically because
prices, entry, and the stationary distribution adjust jointly.

\section{Decision Rules of Exit}

\begin{figure}[htbp]
    \centering
    \includegraphics[width=0.85\textwidth]{../output/task3_exit_policies.pdf}
    \caption{Equilibrium Exit Policies Across Model Specifications}
    \label{fig:exit_policies}
\end{figure}

In the standard model, exit decisions are deterministic: all firms in
the lowest productivity state exit, while firms in higher productivity
states remain. This follows because firms with identical productivity
have identical continuation values and therefore make the same exit
decision.

Introducing TV1 shocks makes exit decisions probabilistic, allowing
firms with identical productivity to make different choices. As a result,
some low-productivity firms survive while some higher-productivity firms
exit. The exit probabilities remain decreasing in productivity, since
more productive firms have higher expected continuation values.
However, the relationship becomes smoother rather than following a
strict cutoff rule.

The parameter $\alpha$ governs the variance of the action-specific
shocks. As $\alpha$ increases, the variance decreases, and firms are
more likely to choose the action with the highest fundamental value.
Consequently, the TV1 exit probabilities approach the deterministic
decision rule as $\alpha$ increases.

\section{Fixed Operating Cost Comparative Statics}

\begin{figure}[htbp]
    \centering
    \includegraphics[width=\textwidth]{../output/task4_fixed_cost_comparison.pdf}
    \caption{Effect of Increasing Fixed Operating Costs from $c_f=10$ to $c_f=15$}
    \label{fig:fixed_cost_comparison}
\end{figure}

\noindent
Increasing the fixed operating cost from $c_f=10$ to $c_f=15$
generally increases the probability of exit, particularly among
low-productivity firms. Higher fixed costs reduce operating profits
and expected continuation values, making continued operation less
attractive. Although the equilibrium output price adjusts through
the free-entry condition, the overall effect in our simulations is
an increase in exit.

In the standard model, the higher fixed cost changes the exit
threshold: firms in the two lowest productivity states now exit,
whereas previously only firms in the lowest state exited. Thus, the
increase in fixed costs induces exit among firms that previously
found continued operation profitable.

In the TV1 specifications, the adjustment occurs through changes
in exit probabilities rather than a discrete change in the exit
threshold. For $\alpha=1$, the exit probabilities in the two lowest
productivity states increase from approximately $(0.52, 0.30)$
to $(0.84, 0.55)$. For $\alpha=2$, they increase from approximately
$(0.67, 0.20)$ to $(0.97, 0.67)$. Exit probabilities for firms in
the three highest productivity states remain close to zero.

These results illustrate that higher operating costs disproportionately
affect firms with low productivity and low continuation values.
Moreover, the TV1 model generates a smoother response to changes
in operating costs, whereas the standard model exhibits discrete
changes in the set of firms choosing to exit.

\end{document}